In 2003 Cohn and Elkies \cite{ElkiesCohn}  developed  linear programming bounds that apply directly to sphere packings. The goal of this section is to formalize the Cohn--Elkies linear programming bound.

The following theorem is the key result of \cite{ElkiesCohn}. (Note that the original theorem is stated for a class of functions more general then Schwartz functions.)
\begin{theorem}[Cohn--Elkies {\cite{ElkiesCohn}}]\label{thm:Cohn-Elkies-periodic}\uses{def:Fourier-Transform, SpherePacking.density}\lean{LinearProgrammingBound'}\leanok

Let $X\subset\mathbb{R}^d$ be a discrete subset such that $\|x-y\|\geq 1$ for any distinct $x,y\in X$. Suppose that $X$ is $\Lambda$-periodic with respect to some lattice $\Lambda\subset\mathbb{R}^d$. Let $f:\R^d\to\R$ be a Schwartz function that is not identically zero and satisfies the following conditions:
\begin{equation}\label{eqn:Cohn-Elkies-condition-1}f(x)\leq 0\mbox{ for } \|x\|\geq 1\end{equation} and
\begin{equation}\label{eqn:Cohn-Elkies-condition-2}\widehat{f}(x)\geq0\mbox{ for all } x\in\R^d.\end{equation}
  Then the density of any $\Lambda$-periodic sphere packing is bounded above by $$\frac{f(0)}{\widehat{f}(0)}\cdot \mathrm{vol}(B_d(0,1/2)).$$
\end{theorem}
\begin{proof}\uses{thm:Poisson-summation-formula}\leanok
Here we reproduce the proof given in \cite{ElkiesCohn}.

The inequality
\begin{equation}\label{eqn: sharp X 1}
\sharp (X/\Lambda)\cdot f(0)\geq \sum_{x\in X}\sum_{y\in X/\Lambda}f(x-y)=\sum_{x\in X/\Lambda}\sum_{y\in X/\Lambda}\sum_{\ell\in  \Lambda}f(x-y+l)\end{equation}
follows from the condition \eqref{eqn:Cohn-Elkies-condition-1} of the theorem and the assumption on the distances between points in $X$.
The equality
$$\sum_{x\in X/\Lambda}\sum_{y\in X/\Lambda}\sum_{\ell\in  \Lambda}f(x-y+l)=\sum_{x\in X/\Lambda}\sum_{y\in X/\Lambda}\frac{1}{\mathrm{vol}(\mathbb{R}^d/\Lambda)}\,\sum_{m\in \Lambda^*} \widehat{f}(m)\,e^{2\pi i m(x-y)}.$$
follows from the Poisson summation formula.
The right hand side of the above equation can be written as
$$\sum_{x\in X/\Lambda}\sum_{y\in X/\Lambda}\frac{1}{\mathrm{vol}(\mathbb{R}^d/\Lambda)}\,\sum_{m\in \Lambda^*} \widehat{f}(m)\,e^{2\pi i m(x-y)}=\frac{1}{\mathrm{vol}(\mathbb{R}^d/\Lambda)}\,\sum_{m\in \Lambda^*} \widehat{f}(m)\cdot\big|\sum_{x\in X/\Lambda}e^{2\pi i m x}\big|^2.$$
Note that $\big|\sum_{x\in X/\Lambda}e^{2\pi i m x}\big|^2\geq0$ for all $m\in\Lambda^*$. Moreover,  the term corresponding to $m=0$ satisfies $\big|\sum_{x\in X/\Lambda}e^{2\pi i 0 x}\big|^2=\sharp (X/\Lambda)^2$.
Now we use the condition \eqref{eqn:Cohn-Elkies-condition-2} and estimate
\begin{equation}\label{eqn: sharp X 2}\frac{1}{\mathrm{vol}(\mathbb{R}^d/\Lambda)}\,\sum_{m\in \Lambda^*} \widehat{f}(m)\cdot\big|\sum_{x\in X/\Lambda}e^{2\pi i m(x-y)}\big|^2
\geq \frac{\sharp (X/\Lambda)^2}{\mathrm{vol}(\mathbb{R}^d/\Lambda)}\cdot \widehat{f}(0).
\end{equation}
Comparing inequalities \eqref{eqn: sharp X 1} and \eqref{eqn: sharp X 2} we arrive at
$$\frac{\sharp (X/\Lambda)}{\mathrm{vol}(\mathbb{R}^d/\Lambda)}\leq \frac{f(0)}{\widehat{f}(0)}.$$
Now we see that the density of the periodic packing $\mathcal{P}_X$ with balls of radius $1/2$ is bounded by
$$\Delta(\mathcal{P}_X)=\frac{\sharp (X/\Lambda)}{\mathrm{vol}(\mathbb{R}^d/\Lambda)}\cdot{\mathrm{vol}(B_d(0,1/2))}\leq
\frac{f(0)}{\widehat{f}(0)}\cdot \mathrm{vol}(B_d(0,1/2)).$$
This finishes the proof of the theorem for periodic packings.
\end{proof}

\begin{theorem}[Cohn--Elkies {\cite{ElkiesCohn}}]\label{thm:Cohn-Elkies-general}\lean{LinearProgrammingBound}\leanok
  Let $f:\R^d\to\R$ be a Schwartz function that is not identically zero and satisfies \eqref{eqn:Cohn-Elkies-condition-1} and \eqref{eqn:Cohn-Elkies-condition-2}. Then the density of any $\Lambda$-periodic sphere packing is bounded above by $$\frac{f(0)}{\widehat{f}(0)}\cdot \mathrm{vol}(B_d(0,1/2)).$$
\end{theorem}
\begin{proof}\uses{thm:Cohn-Elkies-periodic,thm:periodic-packing-optimal}\lean{periodic_constant_eq_constant, periodic_constant_eq_periodic_constant_normalized}\leanok
  The result follows immediately from Theorem~\ref{thm:periodic-packing-optimal} and \Cref{thm:Cohn-Elkies-periodic}.
\end{proof}

\subsection{The Shape of an Optimal Function}

The main step in our proof of \cref{MainTheorem} is the explicit construction of an optimal function for $d = 8$. Before stating what we shall prove, we record what \emph{optimality} forces upon such a function; this analysis, which follows the discussion in \cite[p.~8]{CohnICM}, dictates the entire construction of Section~\ref{sec: fourier double zeroes}.

Suppose $f$ is a Schwartz function satisfying \eqref{eqn:Cohn-Elkies-condition-1} and \eqref{eqn:Cohn-Elkies-condition-2} whose bound in Theorem~\ref{thm:Cohn-Elkies-periodic} is attained by some $\Lambda$-periodic packing with centres $X$. Then every inequality in the proof of Theorem~\ref{thm:Cohn-Elkies-periodic} must be an equality. Equality in \eqref{eqn: sharp X 1} forces $f(x - y) = 0$ for all distinct $x \in X$, $y \in X/\Lambda$; when $X = \Lambda$, this says that \emph{$f$ vanishes at all non-zero lattice points}. Moreover, $f(0) \geq 0$ (else the bound would be negative), so \eqref{eqn:Cohn-Elkies-condition-1} forces these vanishing points of norm $\geq 1$ to be \emph{double} zeroes of the radial profile of $f$: a sign change there would contradict $f \leq 0$ on $\{\|x\| \geq 1\}$. Finally, equality in \eqref{eqn: sharp X 2} forces $\widehat{f}(m) = 0$ for all non-zero $m \in \Lambda^*$; for $\Lambda = \Lambda_8$, which is self-dual, these are the same points again.

Since the conditions \eqref{eqn:Cohn-Elkies-condition-1} and \eqref{eqn:Cohn-Elkies-condition-2} depend only on $\norm{x}$, we may restrict our search to \emph{radial} functions (Definition~\ref{def:radial-function}), and by Lemma~\ref{lemma:eigenfunction-decomposition} it suffices to construct a $+1$- and a $-1$-eigenfunction of the Fourier transform, each with (double) zeroes at the appropriate lattice points, and take a suitable linear combination.

It will be convenient to work at the scale of $\Lambda_8$ itself, whose non-zero vectors have norms $\sqrt{2n}$, $n \geq 1$, rather than rescaling the lattice to separation $1$. We therefore aim for the following.

\begin{theorem}\label{thm:g}\uses{thm:g1}
There exists a radial Schwartz function $g:\R^8\to\R$ which satisfies:
\begin{align}
g(x)&\leq 0\mbox{ for } \|x\|\geq \sqrt{2} \label{eqn:g1}\\
\widehat{g}(x)&\geq0\mbox{ for all } x\in\R^8\label{eqn:g2}\\
g(0)&=\widehat{g}(0)=1.\label{eqn:g3}
\end{align}
\end{theorem}

Theorem \ref{thm:Cohn-Elkies-general} applied to the rescaled function $f(x):=g(\sqrt{2}\, x)$, which satisfies \eqref{eqn:Cohn-Elkies-condition-1} and \eqref{eqn:Cohn-Elkies-condition-2} with $f(0)/\widehat{f}(0) = 16\, g(0)/\widehat{g}(0) = 16$, immediately implies \cref{MainTheorem}: recalling that $\Vol{B_{2k}(0, r)} = r^{2k}\pi^k/k!$ (a fact already available in Mathlib), the resulting bound is
\[
  16 \cdot \Vol{B_8(0, 1/2)} = 16 \cdot \frac{1}{2^8}\cdot\frac{\pi^4}{24} = \frac{\pi^4}{384},
\]
which is precisely the density of the $E_8$ packing. In general the bound is $g(0)/\widehat{g}(0)$ times this density, so it is sharp precisely because of the normalisation \eqref{eqn:g3}: an optimal $g$ must satisfy $g(0) = \widehat{g}(0)$.
